\chapter{Controller design}%
\label{ch:controller_design}

\section{Current controller}

To help fascilitate the design of the position controller, we will first design a controller to control the armature current $I_a$ with the armature voltage $E_a$. This will allow us to treat the armature current as a control input in the position control design, simplifying the dynamics significantly. The current controller will be designed as a PI controller with a feedforward term. Let's first consider the electrical dynamics of the motor from \cref{eq:motor_model_electrical_rewrite}:

\begin{equation}
    \label{eq:controller_design_electrical_dynamics}
    L_a \frac{\mathrm{d}I_a}{\mathrm{d}t} + R_a I_a + K_b \frac{\mathrm{d}\theta}{\mathrm{d}t}= E_a.
\end{equation}

Given the time varying reference current $r(t)$, we can design the input voltage $E_a(t)$ as:

\begin{equation}
    \label{eq:controller_design_voltage_input}
    E_a(t) = \underbrace{K_b \frac{\mathrm{d}\theta(t)}{\mathrm{d}t} + R_a I_a(t)}_{\text{Feedforward term}} + \underbrace{L_a \left[ K_p \left(r(t) - I_a(t)\right) + K_i \int_0^t \left(r(\tau) - I_a(\tau)\right) \, d\tau \right]}_{\text{PI controller}}.
\end{equation}

Substituting \cref{eq:controller_design_voltage_input} into \cref{eq:controller_design_electrical_dynamics}, we get:

\begin{equation}
    \label{eq:controller_design_current_dynamics}
    \frac{\mathrm{d}I_a}{\mathrm{d}t} = K_p (r(t) - I_a(t)) + K_i \int_0^t (r(\tau) - I_a(\tau)) \, d\tau.
\end{equation}

Transforming \cref{eq:controller_design_current_dynamics} to the Laplace domain, we get:

\begin{equation}
    \label{eq:controller_design_current_dynamics_laplace}
    \frac{I_a(s)}{r(s)} = \frac{K_p s + K_i}{s^2 + K_p s + K_i}.
\end{equation}

\textbf{TODO: Compare denominator to the standard second-order system, and provide example values, and plot the current control with noise.}


\section{Controllability of the system}%
\label{sec:controllability}
Based on the state-space representation our system in \cref{subsec:state_space_model_load}, we can check if the system is controllable, as described in \cref{subsec:theory_control_controllability}. The controllability matrix $\mathcal{C}$ for our system is given by:

\begin{equation}
    \label{eq:model_controllability}
    \mathcal{C} = [\mathbf{B}, \mathbf{A}\mathbf{B}, \mathbf{A}^2\mathbf{B}].
\end{equation}

Substituting for our specific $\mathbf{A}$ and $\mathbf{B}$, we get:

\begin{equation}
    \mathcal{C} = 
    \begin{bmatrix}
        0 & 0 & \frac{K_t}{L_a (J_m+J)}\\
        0 & \frac{K_t}{L_a (J_m+J)} & -\frac{K_t [L_a (B_m+c) + R_a (J_m+J)]}{L_a^2 (J_m+J)^2}\\
        \frac{1}{L_a} & -\frac{R_a}{L_a^2} & -\frac{K_b K_t L_a+R_a^2 (J_m+J)}{L_a^3 (J_m+J)}
    \end{bmatrix}.
\end{equation}

Clearly, $\mathrm{rank}(\mathcal{C}) = 3$, which is equal to the number of states. Therefore, the system is controllable.

\section{MRAC implementation}%
\label{sec:mrac_implementation}

